% ============================================================
% esp32_supervisor.tex — low-voltage side, Path B target.
%
% Board-level view of the enclosure interior: isolated 5 V rail
% feeding the ESP32, MAX31855 SPI thermocouple front-end, NTC
% frame-temperature divider, opto-isolated relay module driving
% the contactor coil circuit, local ack button, onboard LED.
%
% Nothing left of the dashed boundary sits at mains potential.
%
% Layout is banded to keep every conductor orthogonal and
% crossing-free:
%   y = 7.0 .. 9.0   SPI bus + MAX31855
%   y = 2.6 .. 5.4   NTC frame divider
%   y = -0.7 .. 3.0  relay module + pulldown
%   x < 0            power in, ack, status LED
%
% Regenerate with:  make esp32_supervisor.svg
% ============================================================
\def\pgfsysdriver{pgfsys-dvisvgm.def}
\documentclass[border=10pt,dvisvgm]{standalone}
% ============================================================
% tsstyle.tex — shared preamble for every sheet in this directory.
%
% Inputted by each sheet after \documentclass. Sheets stay pure
% geometry; everything about how the geometry *looks* lives here,
% so a style change is a one-file edit.
% ============================================================

\usepackage[T1]{fontenc}
\usepackage[utf8]{inputenc}
\usepackage[scaled=0.92]{helvet}
\renewcommand{\familydefault}{\sfdefault}

\usepackage[american,siunitx,RPvoltages]{circuitikz}
\usetikzlibrary{calc,positioning,fit,backgrounds,arrows.meta,decorations.pathreplacing}

% ------------------------------------------------------------
% Palette. Deliberately small: ink for conductors, mute for
% annotation, and exactly two semantic colours — mains (danger)
% and SELV (safe). Colour carries meaning on these sheets; it is
% never decoration.
% ------------------------------------------------------------
\definecolor{tsink}{HTML}{16181D}    % conductors, symbol bodies
\definecolor{tsmute}{HTML}{6E7178}   % secondary annotation
\definecolor{tsrule}{HTML}{C7CAD1}   % boxes, leader lines
\definecolor{tsmains}{HTML}{B3261E}  % mains potential — danger
\definecolor{tsselv}{HTML}{15607A}   % SELV / low-voltage rail
\definecolor{tssafe}{HTML}{1E6B45}   % fail-safe / passive layer
\definecolor{tspanel}{HTML}{F5F6F7}  % block fills
\definecolor{tspanelb}{HTML}{E4E7EB} % block fills, alternate

% ------------------------------------------------------------
% Global circuitikz sizing. Bipoles are scaled down from the
% default because these sheets carry a lot of elements per unit
% area; the default American resistor is comically large next to
% an IC block.
% ------------------------------------------------------------
\ctikzset{
  bipoles/length=0.85cm,
  resistors/scale=0.85,
  capacitors/scale=0.8,
  label/align=straight,
  font=\footnotesize,
}

\tikzset{
  % Conductors -------------------------------------------------
  wire/.style        = {tsink, line width=0.9pt, line cap=round},
  buswire/.style     = {tsink, line width=1.5pt, line cap=round},
  mainswire/.style   = {tsmains, line width=1.5pt, line cap=round},
  selvwire/.style    = {tsselv, line width=1.2pt, line cap=round},
  % Functional blocks -----------------------------------------
  block/.style       = {draw=tsink, line width=1pt, fill=tspanel,
                        rounded corners=1.5pt, align=center,
                        inner sep=6pt, font=\small\bfseries},
  subblock/.style    = {draw=tsrule, line width=0.7pt, fill=white,
                        rounded corners=1.5pt, align=center,
                        inner sep=4pt, font=\scriptsize},
  % Text ------------------------------------------------------
  pinlbl/.style      = {font=\scriptsize, tsink, inner sep=2pt},
  siglbl/.style      = {font=\scriptsize\itshape, tsmute, inner sep=1.5pt},
  netlbl/.style      = {font=\scriptsize\bfseries, tsink, inner sep=2pt},
  note/.style        = {font=\scriptsize, tsmute, align=left},
  danger/.style      = {font=\scriptsize\bfseries, tsmains, align=left},
  % Isolation boundary ----------------------------------------
  isoline/.style     = {tsmains, line width=1pt, dash pattern=on 5pt off 3pt},
  % A wire that crosses another without connecting. The white
  % under-stroke reads as an over-pass, so a crossing can never be
  % mistaken for a junction (junctions are always a filled dot).
  overpass/.style    = {preaction={draw, white, line width=3.4pt}},
  % Sheet furniture -------------------------------------------
  titleblock/.style  = {draw=tsink, line width=1pt, fill=white,
                        align=left, inner sep=7pt, font=\scriptsize},
  legendbox/.style   = {draw=tsrule, line width=0.7pt, fill=white,
                        align=left, inner sep=7pt, font=\scriptsize},
}

% ------------------------------------------------------------
% \pinrow{<x>}{<y>}{<dir>}{<label>}  — a pin stub on a block
% edge. dir is -1 (left edge, stub runs left) or +1 (right edge).
% ------------------------------------------------------------
\newcommand{\pinstub}[5]{%
  \draw[wire] (#1,#2) -- ++(#3*0.55,0);
  \node[pinlbl, anchor=#4] at (#1+#3*0.10,#2+0.14) {#5};
}

% ------------------------------------------------------------
% \titleblocknode{<x>}{<y>}{<sheet>}{<subtitle>}{<source>}
% Standard title block, anchored at its south east corner.
% ------------------------------------------------------------
\newcommand{\sheettitle}[5]{%
  \node[titleblock, anchor=south east] at (#1,#2) {%
    {\normalsize\bfseries #3}\\[2pt]
    {\color{tsmute}#4}\\[4pt]
    {\color{tsmute}Table-saw thermal protection retrofit\hfill}\\
    {\color{tsmute}Generated from \texttt{#5} — do not edit the SVG}%
  };
}


\begin{document}
\begin{circuitikz}[line width=0.9pt]

% ============================================================
% Provenance banner — which hardware this sheet is drawn for
% ============================================================
\node[font=\scriptsize\bfseries, tsmains, anchor=south west] at (-4.6,10.6)
  {PATH B TARGET --- MAX31855, RELAY MODULE AND ISOLATED PSU ARE
   \textbf{NOT PURCHASED} (TASK-1/2/3).};
\node[font=\scriptsize, tsmute, anchor=south west] at (-4.6,10.2)
  {For the hardware actually on hand, build from
   \texttt{pathA\_supervisor.svg} and \texttt{WIRING-PATH-A.md}.};

% ============================================================
% ESP32-DevKitC-32E
%
% The purchased board (B0GF1ZJCCN) is the 38-pin USB-C variant.
% Every GPIO used below is present on it. The pin assignments
% are unchanged from the 30-pin DOIT-style board, which also
% carries all of them — but the label has to match the part.
% ============================================================
\draw[block, fill=tspanelb] (0,0.6) rectangle (4.2,9.8);
\node[font=\small\bfseries, tsink, align=center] at (2.1,5.6)
  {ESP32-\\DevKitC-32E};
\node[font=\scriptsize, tsmute, align=center] at (2.1,4.6)
  {\texttt{B0GF1ZJCCN}\\38-pin, USB-C\\3.3\,V logic};

% ---- left edge: power ----
\node[pinlbl, anchor=west] at (0.12,9.0) {VIN};
\node[pinlbl, anchor=west] at (0.12,8.0) {3V3};
\node[pinlbl, anchor=west] at (0.12,7.0) {GND};
\node[pinlbl, anchor=west] at (0.12,6.0) {GND};

% +5 V in from the isolated PSU
\draw[selvwire] (0,9.0) -- (-3.0,9.0);
\node[netlbl, anchor=south, tsselv] at (-1.5,9.08) {+5\,V};
\node[note, anchor=east, tsselv] at (-3.1,9.0) {isolated PSU\\(SELV)};
\draw[selvwire] (-3.0,9.15) -- (-3.0,8.85);

% 3V3 rail — regulated on the module, feeds MAX31855 + divider
\draw[wire] (0,8.0) -- (-2.4,8.0);
\node[netlbl, anchor=south] at (-1.2,8.08) {3V3};
\node[note, anchor=east] at (-2.5,8.0) {regulated on\\the DevKitC};

% Both GND pins tied to one common bus — see WIRING.md checklist #1
\draw[wire] (0,7.0) -- (-1.3,7.0);
\draw[wire] (0,6.0) -- (-1.3,6.0);
\draw[wire] (-1.3,7.0) -- (-1.3,5.4);
\draw (-1.3,6.0) node[circ]{};
\draw (-1.3,5.4) node[ground]{};
\node[note, anchor=east] at (-1.6,6.6) {common GND bus\\(both pins jumpered)};

% ---- left edge: operator I/O ----
\node[pinlbl, anchor=west] at (0.12,2.6) {GPIO2};
\node[pinlbl, anchor=west] at (0.12,1.4) {GPIO27};

% Status LED — onboard on the DevKitC, no external wiring
\draw[wire] (0,2.6) -- (-1.4,2.6);
\draw (-1.4,2.6) node[ocirc]{};
\node[note, anchor=east] at (-1.65,2.6) {onboard blue LED\\(internal to module)};

% Ack button — GPIO27 to GND, firmware enables the internal pullup
\draw[wire] (0,1.4) -- (-3.4,1.4);
\draw[wire] (-3.4,1.4) to[push button] (-3.4,0.0);
\draw (-3.4,0.0) node[ground]{};
\node[note, anchor=east] at (-3.75,0.85) {ACK button\\(momentary NO,\\firmware pullup)};

% ---- right edge: signals ----
\node[pinlbl, anchor=east] at (4.08,9.0) {GPIO18};
\node[pinlbl, anchor=east] at (4.08,8.0) {GPIO19};
\node[pinlbl, anchor=east] at (4.08,7.0) {GPIO5};
\node[pinlbl, anchor=east] at (4.08,3.8) {GPIO34};
\node[pinlbl, anchor=east] at (4.08,1.4) {GPIO26};

% ============================================================
% SPI bus — three straight runs, no crossings
% ============================================================
\draw[wire] (4.2,9.0) -- (9.6,9.0);
\draw[wire] (4.2,8.0) -- (9.6,8.0);
\draw[wire] (4.2,7.0) -- (9.6,7.0);
\node[siglbl, anchor=south] at (6.9,9.08) {SCK};
\node[siglbl, anchor=south] at (6.9,8.08) {MISO};
\node[siglbl, anchor=south] at (6.9,7.08) {CS (active-LOW)};

% ============================================================
% MAX31855 — K-type thermocouple front-end
% ============================================================
\draw[block] (9.6,6.2) rectangle (13.2,9.8);
\node[font=\small\bfseries, tsink, align=center] at (11.4,7.5)
  {MAX31855};
\node[font=\scriptsize, tsmute, align=center] at (11.4,6.75)
  {K-type, SPI\\3.3\,V only};
\node[pinlbl, anchor=west] at (9.72,9.0) {SCK};
\node[pinlbl, anchor=west] at (9.72,8.0) {SO};
\node[pinlbl, anchor=west] at (9.72,7.0) {CS};
\node[pinlbl, anchor=east] at (13.08,8.8) {T+};
\node[pinlbl, anchor=east] at (13.08,7.6) {T$-$};

% VCC — 3.3 V, never 5 V
\draw[wire] (11.0,9.8) -- (11.0,10.5);
\draw (11.0,10.5) node[vcc]{3V3};
\node[note, anchor=west] at (11.25,10.25) {\textbf{not} 5\,V};

% GND
\draw[wire] (11.0,6.2) -- (11.0,5.5);
\draw (11.0,5.5) node[ground]{};

% Thermocouple leads out to the motor pigtail
\draw[wire] (13.2,8.8) -- (15.4,8.8);
\draw[wire] (13.2,7.6) -- (15.4,7.6);
\draw (15.4,8.8) node[ocirc]{};
\draw (15.4,7.6) node[ocirc]{};
\node[note, anchor=west] at (15.65,8.8) {TC+ \textbf{yellow} $\rightarrow$ winding};
\node[note, anchor=west] at (15.65,7.6) {TC$-$ \textbf{red} $\rightarrow$ winding};
\node[note, anchor=west, tssafe] at (15.65,6.75)
  {K-type is polarised.\\Red is the \emph{negative} lead.\\AlN substrate at the tip\\isolates from winding potential.};

% ============================================================
% NTC frame divider:  3V3 -> 10k -> GPIO34 -> NTC -> GND
% ============================================================
\draw[wire] (4.2,3.8) -- (6.8,3.8);
\node[siglbl, anchor=south] at (5.5,3.88) {ADC1\_CH6};
\draw[wire] (6.8,3.8) to[R] (6.8,5.1);
\draw[wire] (6.8,5.1) -- (6.8,5.4);
\draw (6.8,5.4) node[vcc]{3V3};
\node[pinlbl, anchor=west] at (7.05,4.45) {10\,k$\Omega$ 1\%};

\draw[wire] (6.8,3.8) to[thermistor] (6.8,2.9);
\draw[wire] (6.8,2.9) -- (6.8,2.6);
\draw (6.8,2.6) node[ground]{};
\draw (6.8,3.8) node[circ]{};
\node[pinlbl, anchor=west] at (7.05,3.25) {NTC 10\,k$\Omega$};
\node[note, anchor=west] at (7.05,2.85) {B3950, on frame\\(advisory only)};

% ============================================================
% Relay module — opto-isolated, active-HIGH, NO contacts
% ============================================================
\draw[wire] (4.2,1.4) -- (9.6,1.4);
\node[siglbl, anchor=south] at (5.6,1.48) {relay drive};

% 10k pulldown — non-optional. Floating GPIO26 must leave the relay OPEN.
\draw (8.4,1.4) node[circ]{};
\draw[wire] (8.4,1.4) to[R] (8.4,0.1);
\draw[wire] (8.4,0.1) -- (8.4,-0.2);
\draw (8.4,-0.2) node[ground]{};
\node[pinlbl, anchor=east] at (8.15,0.78) {10\,k$\Omega$};
\node[note, anchor=east, tssafe] at (8.15,0.38) {pulldown\\(mandatory)};

\draw[block] (9.6,0.0) rectangle (13.6,3.0);
\node[font=\small\bfseries, tsink, align=center] at (11.6,1.85)
  {Relay module};
\node[font=\scriptsize, tsmute, align=center] at (11.6,1.0)
  {opto-isolated\\active-HIGH, N.O.\\5\,V coil};
\node[pinlbl, anchor=west] at (9.72,1.4) {IN};
\node[pinlbl, anchor=east] at (13.48,2.3) {COM};
\node[pinlbl, anchor=east] at (13.48,1.4) {NO};
\node[pinlbl, anchor=east] at (13.48,0.5) {NC};

\draw[selvwire] (10.6,3.0) -- (10.6,3.7);
\draw (10.6,3.7) node[vcc, color=tsselv]{+5\,V};
\draw[wire] (10.6,0.0) -- (10.6,-0.7);
\draw (10.6,-0.7) node[ground]{};

% ---- contacts out to the mains-side coil circuit ----
\draw[mainswire] (13.6,2.3) -- (15.4,2.3);
\draw[mainswire] (13.6,1.4) -- (15.4,1.4);
\draw[wire] (13.6,0.5) -- (14.4,0.5);
\draw (15.4,2.3) node[ocirc, color=tsmains]{};
\draw (15.4,1.4) node[ocirc, color=tsmains]{};
\draw (14.4,0.5) node[ocirc]{};
\node[note, anchor=west, tsmains] at (15.65,2.3)
  {$\rightarrow$ thermostat NC lead};
\node[note, anchor=west, tsmains] at (15.65,1.4)
  {$\rightarrow$ M1 coil terminal A1};
\node[note, anchor=north] at (14.1,0.34) {leave open};

% ---- isolation boundary ----
\draw[isoline] (14.9,-0.9) -- (14.9,3.2);
\node[danger, anchor=south west] at (14.55,3.32)
  {MAINS POTENTIAL BEYOND THIS LINE};

% ============================================================
% Isolation legend
% ============================================================
\node[legendbox, anchor=north west, text width=7.3cm] at (-4.6,-1.6) {%
  \textbf{Three isolation barriers} (all external to this sheet)\\[3pt]
  \begin{tabular}{@{}p{1.85cm}@{\hspace{4pt}}p{4.9cm}@{}}
    PSU input   & reinforced, UL 61010 / IEC 60950 rated part\\[1pt]
    Thermocouple& AlN dielectric substrate at the winding\\[1pt]
    Relay       & opto-input \emph{and} contact-gap output\\
  \end{tabular}
};

\node[legendbox, anchor=north west, text width=7.6cm] at (3.4,-1.6) {%
  \textbf{Fail-safe behaviour}\\[3pt]
  \textcolor{tssafe}{\textbf{SR-1}}\quad Relay is N.O., energised-to-close.
  Floating GPIO26 $=$ relay open $=$ coil dropped.\\[2pt]
  \textcolor{tssafe}{\textbf{SR-3}}\quad The passive thermostat sits in series
  ahead of COM, on the mains side. It cuts the coil with no
  firmware involvement.\\[2pt]
  Coil circuit is \emph{not} on this sheet — see
  \texttt{ladder\_coil\_circuit.svg}.
};

% ============================================================
% Title block
% ============================================================
\node[titleblock, anchor=north east, text width=6.2cm] at (22.6,-1.6) {%
  {\normalsize\bfseries Low-voltage supervisor}\\[2pt]
  {\color{tsmute}Signal schematic — design of record}\\[5pt]
  {\color{tsmute}Table-saw thermal protection retrofit}\\
  {\color{tsmute}Pin-level truth: \texttt{WIRING.md}}\\
  {\color{tsmute}Source: \texttt{esp32\_supervisor.tex}}%
};

\end{circuitikz}
\end{document}
